\documentclass[10pt,a4paper]{amsart}
\usepackage[margin=19mm]{geometry}
\usepackage{amsmath,amssymb,mathtools}
\usepackage[scr=rsfs,cal=euler]{mathalfa}
\usepackage{xfrac}
\usepackage[unicode,hidelinks,bookmarks=false]{hyperref}
\newcommand{\ie}{i.e.\ }
\newcommand{\cf}{cf.\ }
\newcommand{\resp}{respectively\ }
\newcommand{\Subsection}{\subsection}
\providecommand{\nobreakdash}{\nobreak\hbox{-}}
\setlength{\emergencystretch}{2em}
\multlinegap0pt
\usepackage{bm}
\usepackage{bbm}
\usepackage{dsfont}
\usepackage{xspace}
\usepackage{cancel}
\usepackage{mathtools}
\usepackage{amsthm}
\makeatletter
\@ifundefined{theorem}{\newtheorem{theorem}{Theorem}}{}
\@ifundefined{lemma}{\newtheorem{lemma}[theorem]{Lemma}}{}
\makeatother
\newcommand{\D}{\mathcal{D}}
\newcommand{\G}{\mathcal{G}}
\newcommand{\K}{\mathcal{K}}
\newcommand{\abs}[1]{\lvert#1\rvert}
\newcommand{\h}{\mathcal{H}}
\newcommand{\norm}[1]{\lVert#1\rVert}
\title[Normalized solutions of nonlinear Dirac equations on noncomp]{Normalized solutions of nonlinear Dirac equations on noncompact metric graphs with localized nonlinearities}
\author{Zhentao He, Chao Ji}
\date{Source: arXiv:2505.15100v1 (CC BY 4.0)}
\begin{document}
\maketitle

\noindent\emph{Source and licence.} Normalized solutions of nonlinear Dirac equations on noncompact metric graphs with localized nonlinearities. Version arXiv:2505.15100v1 (CC BY 4.0), distributed under the Creative Commons Attribution 4.0 International licence, as recorded in the arXiv licence metadata for this exact version and shown on the abstract page https://arxiv.org/abs/2505.15100v1. Full rights notice: licenses/arxiv-2505.15100-CC-BY-4.0.txt.\par\smallskip

\noindent\textit{Editorial scope.} Counted unit: the Lemma whose source label is \texttt{lemcin} in the article (source0.tex lines 757--759, source label \texttt{lemcin}), delivered together with the article's own complete original proof of it (source0.tex lines 760--825, one \texttt{proof} environment of 66 lines). No mathematics is written, repaired or re-derived: statement and proof are the article's own text line for line. Required context retained uncounted: lemgnsg (lines 635--644); eqcssup (lines 751--753). Every definition, display and label of those lines is retained unchanged. Omitted from this package and not redistributed: 1--634, 645--750, 754--756, 826--1864. Nothing inside a retained excerpt is omitted or altered: every retained body is delivered exactly as the article's lines are written, so no omission receipt of any kind accompanies this package. Citations: the retained excerpts carry no citation command at all, so no bibliography entry of the article is transcribed and no reference is redistributed. Reference adaptation: none was needed and none was made; every reference inside the delivered unit resolves inside it, so no unresolved reference or citation is produced. Printed slips kept literal: no printed word, symbol, hypothesis or number of the counted statement or of its proof was altered. Layout and class: the article's own class is \texttt{amsart} with packages including amsmath, amssymb, amsthm, latexsym, cite, cancel, caption, graphicx, wasysym, overpic, tikz, color, subfigure, hyperref; the wrapper is the lane's pinned \texttt{amsart} editorial wrapper at 10pt on A4, which restates the theorem-like environments the retained text needs and adds the article's own macro definitions that the retained text needs (\texttt{\textbackslash D}, \texttt{\textbackslash G}, \texttt{\textbackslash K}, \texttt{\textbackslash abs}, \texttt{\textbackslash h}, \texttt{\textbackslash norm}), with extra wrapper packages (bm, bbm, dsfont, xspace, cancel, mathtools). The delivered numbering is the unit's own. Assets: no retained span contains a figure environment, a graphics-inclusion command, a PDF-page inclusion, a table or a TikZ picture; no image file is copied into this package, the package declares no asset dependency and no image file is redistributed with it. Licence: version arXiv:2505.15100v1 of this article is distributed under the Creative Commons Attribution 4.0 International licence, as recorded in the arXiv licence metadata for this exact version and shown on the abstract page https://arxiv.org/abs/2505.15100v1. A full rights notice is delivered at licenses/arxiv-2505.15100-CC-BY-4.0.txt. Characters: every retained excerpt is pure ASCII, so the delivered engine, XeLaTeX with its default Latin Modern text font, reports no missing character.
\begin{lemma}\label{lemgnsg}
Let $p \geq 2$, there exist constants $C_{p,\K}$ and $C_{p,\G}$ that depend only on $c$, $m$, $p$ and $\G$ such that
\[
    \int_{\K} \abs{u}^p\,dx \leq C_{p,\K} \norm{u}^{p-2}\norm{u}_2^{2}, \quad \forall u \in Y,
\]
and
\[
 \int_{\K} \abs{u}^p\,dx \leq \int_{\G} \abs{u}^p\,dx \leq C_{p,\G} \norm{u}^{p-2}\norm{u}_2^{2}, \quad \forall u \in Y.
\]
\end{lemma}

\begin{equation}\label{eqcssup}
    c_{\infty} \leq \sup _{0 \leq t<|v|_{2}^{-1}} J(t v).
\end{equation}

\begin{lemma}\label{lemcin}
    Let $2<p<4$. Then $c_\infty < \frac{m}{2}$.
\end{lemma}

\begin{proof}
Let the half-lines of the graph $\G$ be denoted by $\h_1, \h_2,...,\h_N$, and let the total length of the graph $\K$ be $\abs{\K} = \sum_{e \in \K}\ell_e$. Let $b > 0$. Define $\varphi_{b}^1: \G \to \mathbb{C}$ as
$$
    \varphi_{b}^1(x)= \begin{cases}1 & \text { for } x \in \K\\ \max\{0,1-bx\} & \text { for } x \in \h_i,\quad i = 1,2,...,N, \end{cases}
$$
and define $\varphi_{b}: \G \to \mathbb{C}^2$ by
$$\varphi_{b} = \binom{\varphi_{b}^1}{0}.$$

We first estimate  $J(\frac{\varphi_b^+}{\norm{\varphi_b^+}_2})$.
It is clear that $\varphi_{b} \in \operatorname{dom}(\D)$, $(\D- mI)\varphi_{b} = \left(0, -\imath \left(\varphi_{b}^1\right)^\prime\right)^T$ and $((\D- mI)\varphi_{b}, \varphi_{b})_2 = 0$. Thus, we obtain
\[
\norm{\left(\varphi_{b}^1\right)^\prime}_{L^2(\G)}\norm{\varphi_{b}^-}_2\geq \abs{((\D- mI)\varphi_{b},\varphi_{b}^-)_2} = \abs{((\D- mI)\varphi_{b}^-,\varphi_{b}^-)_2} = \norm{\varphi_{b}^-}^2 + m \norm{\varphi_{b}^-}^2_2 \geq 2m\norm{\varphi_{b}^-}_2^2.
\]
Hence,
\begin{equation}\label{eqphi-}
\norm{\varphi_{b}^-}^2 + m \norm{\varphi_{b}^-}^2_2 \leq \abs{((\D- mI)\varphi_{b},\varphi_{b}^-)_2} \leq \frac{1}{2m}\norm{\left(\varphi_{b}^1\right)^\prime}_{L^2(\G)}^2 = \frac{Nb}{2m}.
\end{equation}
Since $((\D- mI)\varphi_{b},\varphi_{b}^+)_2 + ((\D- mI)\varphi_{b},\varphi_{b}^-)_2 = ((\D- mI)\varphi_{b}, \varphi_{b})_2 = 0$, we obtain
\begin{equation}\label{eqphi+}
\norm{\varphi_{b}^+}^2 - m \norm{\varphi_{b}^+}^2_2 = ((\D- mI)\varphi_{b}^+,\varphi_{b}^+)_2 = ((\D- mI)\varphi_{b},\varphi_{b}^+)_2 \leq \frac{Nb}{2m}.
\end{equation}
Since $\norm{\varphi_b}_2^2 = \frac{N}{3b} + \abs{\K}$, \eqref{eqphi-} implies
\begin{equation}\label{eqphi+2}
  \frac{\norm{\varphi_b^+}^2_2}{\norm{\varphi_b}^2_2} \to 1, \quad \frac{\norm{\varphi_b^+}^2}{\norm{\varphi_b}^2} \to 1 \text{ and } \int_{\K}\abs{\varphi_b^+}^{p} \to \int_{\K}\abs{\varphi_b}^{p} = \abs{\K} \text{ as } b \to 0.
\end{equation}

Next, for $t > 0$, from the definition of $h$, we have
$$
I_0(t\varphi_b+h(t\varphi_b)) \geq I_0(t\varphi_b),
$$
which implies
\begin{equation}\label{eqhphi}
    \frac{p}{2}\norm{h(t\varphi_b)}^{2} \leq a\int_{\K}\abs{t\varphi_b+h(t\varphi_b)}^{p}\,dx + \frac{p}{2}\norm{h(t\varphi_b)}^{2} \leq a\int_{\K}\abs{t\varphi_b}^{p} \,dx = a t^p\abs{\K}.
\end{equation}


Thus, for sufficiently small $b>0$, using \eqref{eqphi+}-\eqref{eqhphi}, $p \in (2,4)$, Gagliardo-Nirenberg-Sobolev inequality in Lemma \ref{lemgnsg} and convexity of $\Psi$, we obtain
    \begin{equation}\label{eqjsmallm}
            \begin{aligned}
                J(\frac{\varphi_b^+}{\norm{\varphi_b^+}_2}) =& \frac{\norm{\varphi_b^+}^2}{2\norm{\varphi_b^+}^2_2}-\frac{1}{2}\left\lVert h\left(\frac{\varphi_b^+}{\norm{\varphi_b^+}_2}\right) \right\rVert^2 - \Psi\left(\frac{\varphi_b^+}{\norm{\varphi_b^+}_2}+h\left(\frac{\varphi_b^+}{\norm{\varphi_b^+}_2}\right)\right)\\
                \leq& \frac{m}{2} + \frac{\norm{\varphi_b^+}^2 - m\norm{\varphi_b^+}^2_2}{2\norm{\varphi_b^+}^2_2}   -\frac{1}{2}\left\lVert h\left(\frac{\varphi_b^+}{\norm{\varphi_b^+}_2}\right) \right\rVert^2\\
                & - 2^{1-p}\norm{\varphi_b^+}_2^{-p}\Psi(\varphi_b^+) +  \Psi\left(h\left(\frac{\varphi_b^+}{\norm{\varphi_b^+}_2}\right)\right)\\
                \leq& \frac{m}{2} + C_1b^2 - C_2ab^\frac{p}{2} - \frac{1}{2}\left\lVert h\left(\frac{\varphi_b^+}{\norm{\varphi_b^+}_2}\right)\right\rVert^2\left(1 - \frac{aC_{p,\K}}{pm}\left\lVert h\left(\frac{\varphi_b^+}{\norm{\varphi_b^+}_2}\right) \right\rVert^{p-2}\right) \\
                \leq& \frac{m}{2} + C_1b^2 - C_2ab^\frac{p}{2} - \frac{1}{2}\left\lVert h\left(\frac{\varphi_b^+}{\norm{\varphi_b^+}_2}\right)\right\rVert^2(1 - C_3 a^\frac{p}{2}b^\frac{p(p-2)}{4})\\
                <& \frac{m}{2}.
    \end{aligned}
    \end{equation}

Define $g(t):=J(t\varphi_b^+)$. Then, from \eqref{eqhphi} and the definition of $J$, for $0 \leq t \leq\norm{\varphi_b^+}_{2} ^{-1}$ and sufficiently small $b$, we have
$$
\begin{aligned}
g'(t) & =\frac{1}{t}\langle J^{\prime}(t\varphi_b^+), t\varphi_b^+\rangle=\frac{1}{t}\langle I_0'(t\varphi_b^++h(t\varphi_b^+)), t \varphi_b^++h'(t\varphi_b^+) t\varphi_b^+\rangle \\
& =\frac{1}{t}\langle I_0^{\prime}(t\varphi_b^++h(t\varphi_b^+)), t\varphi_b^+ + h(t\varphi_b^+)\rangle \\
& =\frac{1}{t}\left[\norm{t\varphi_b^+}^2-\norm{h(t\varphi_b^+)}^2-a\int_{\K} \abs{t\varphi_b^+ + h(t\varphi_b^+)}^p \,dx\right]\\
& \geq \frac{1}{t}\left[\norm{t\varphi_b^+}^{2}-a\int_{\K}\abs{t\varphi_b^+}^{p}\,dx+\frac{p-2}{2}\norm{h(t\varphi_b^+)}^{2}\right] \\
& \geq t\norm{\varphi_b^+}^2-at^{p-1}\int_{\K} \abs{\varphi_b^+}^{p} \,dx \\
& \geq t\left(\norm{\varphi_b^+}^2-a\norm{\varphi_b^+}_{2}^{2-p}\int_{\K} \abs{\varphi_b^+}^{p} \,dx\right)\\
& \geq t(C_1b^{-1} - C_2ab^\frac{p-2}{2}),
\end{aligned}
$$
which implies $g(t)$ is increasing for $0 \leq t \leq\norm{\varphi_b^+}_{2} ^{-1}$ when $b$ is sufficiently small. Now, fix $b$ sufficiently small. Then, by \eqref{eqcssup} and \eqref{eqjsmallm}, we have
$$
c_{\infty} \leq \sup _{0 \leq t \leq\norm{\varphi_b^+}_{2} ^{-1}} J(t\varphi_b^+) \leq J(\frac{\varphi_b^+}{\norm{\varphi_b^+}_2}) < \frac{m}{2}.
$$
This completes the proof of lemma.
\end{proof}

\end{document}
